% !TeX root = ../main.tex
\documentclass[../main.tex]{subfiles}

\begin{document}

\section{Diseño del sistema}

\subsection{Requisitos y decisiones de diseño del sistema}

En esta sección se definen los requisitos que debe cumplir el sistema de sonificación EEG--MIDI desarrollado en este Trabajo Fin de Grado. Estos requisitos permiten delimitar qué debe hacer el prototipo, bajo qué condiciones debe operar y qué decisiones de diseño se han adoptado en la implementación final.

El sistema se plantea como un prototipo experimental de sonificación EEG--MIDI en vivo. Su finalidad no es realizar diagnóstico médico ni evaluar patologías neurológicas, sino transformar señales EEG reales en eventos musicales enviados mediante una salida MIDI física.

\subsubsection{Requisitos funcionales}

Los requisitos funcionales definen las capacidades principales que debe ofrecer el sistema. En este proyecto abarcan la cadena completa de funcionamiento: adquisición EEG real, transmisión de muestras, procesamiento, extracción de características, generación musical y salida MIDI física.

\begin{description}[leftmargin=!,labelwidth=1.8cm]

    \item[RF-1] \textbf{Adquisición EEG en vivo.} El sistema debe permitir la adquisición de señales EEG reales procedentes de un usuario durante una sesión de funcionamiento. La señal utilizada para la sonificación debe proceder de una captura en vivo y no de la reproducción \textit{offline} de bases de datos o registros EEG previamente almacenados.

    Este requisito responde al objetivo principal del proyecto: transformar actividad EEG real en eventos musicales MIDI durante el funcionamiento del prototipo.

    \item[RF-2] \textbf{Adquisición multicanal de señales EEG.} El sistema debe permitir la adquisición de señales EEG en varios canales de forma simultánea. Aunque la sonificación inicial se base en un canal principal, la adquisición multicanal permite conservar información adicional y facilita futuras ampliaciones.

    La implementación desarrollada utiliza el ADS1299--4PAG, pero el requisito se centra en la capacidad multicanal del sistema, no en cerrar la arquitectura a un único componente.

    \item[RF-3] \textbf{Configuración flexible de adquisición.} El sistema debe permitir configurar parámetros de adquisición, especialmente la frecuencia de muestreo, para adaptarse a distintas necesidades de análisis EEG.

    La frecuencia de \SI{250}{\hertz} utilizada en las pruebas del TFG se considera una decisión de diseño, no una limitación funcional del sistema.

    \item[RF-4] \textbf{Compatibilidad con distintos montajes EEG de superficie.} El sistema debe permitir la conexión de electrodos en diferentes configuraciones EEG de superficie, incluyendo posiciones compatibles con el sistema internacional 10--20.

    La configuración ear--EEG/mastoides con referencia en muñeca izquierda corresponde al montaje empleado en la validación, no a una restricción funcional obligatoria.

    \item[RF-5] \textbf{Transmisión continua de muestras.} El sistema debe transmitir de forma continua las muestras EEG digitalizadas desde el bloque de adquisición hacia el bloque de procesamiento.

    Este requisito garantiza que la señal capturada pueda analizarse durante la sesión activa y emplearse posteriormente en la generación musical.

    \item[RF-6] \textbf{Procesamiento continuo de la señal EEG.} El sistema debe procesar de forma continua las muestras EEG recibidas, organizándolas en buffers o ventanas temporales.

    La implementación actual se realiza en Python sobre Arduino UNO Q, pero el requisito funcional es el procesamiento continuo, no el lenguaje empleado.

    \item[RF-7] \textbf{Extracción de características EEG.} El sistema debe extraer características temporales y espectrales de la señal EEG adquirida, como amplitud, RMS, potencia absoluta, potencia relativa y frecuencias pico.

    Estas características actúan como representación intermedia entre la señal EEG cruda y el módulo de sonificación.

    \item[RF-8] \textbf{Análisis por bandas EEG.} El sistema debe obtener información asociada a bandas EEG de interés, como delta, theta, alfa, beta y gamma.

    Estas bandas permiten describir el contenido frecuencial de la señal y utilizarlo como entrada para el mapeo musical.

    \item[RF-9] \textbf{Generación de parámetros musicales a partir de características EEG.} El sistema debe transformar las características EEG extraídas en parámetros musicales como altura, duración, intensidad, escala, tonalidad o densidad rítmica.

    Este requisito constituye el núcleo de la sonificación EEG--MIDI, al convertir información cerebral en decisiones musicales.

    \item[RF-10] \textbf{Organización musical jerárquica.} El sistema debe organizar la generación musical mediante una estructura jerárquica formada por segmentos, compases y notas.

    Esta organización evita una conversión directa de señal a sonido y permite generar una salida musical más estructurada.

    \item[RF-11] \textbf{Generación de eventos MIDI válidos.} El sistema debe generar eventos MIDI válidos, especialmente mensajes equivalentes a \textit{Note On} y \textit{Note Off}.

    Cada nota debe tener inicio, duración y finalización definidos para evitar mensajes incompletos o notas colgadas.

    \item[RF-12] \textbf{Salida MIDI física.} El sistema debe enviar los eventos musicales generados mediante una salida MIDI física.

    En la implementación desarrollada se emplea la línea asociada al pin D1/RX y la PCB MIDI, pero el requisito principal es disponer de una salida física operativa.

    \item[RF-13] \textbf{Monitorización auxiliar del sistema.} El sistema debe disponer de una interfaz auxiliar de monitorización que permita observar el estado general del prototipo durante su funcionamiento.

    La interfaz web actúa como apoyo para supervisar adquisición, procesamiento y generación musical, sin sustituir a la salida MIDI física.

    \item[RF-14] \textbf{Separación entre adquisición, procesamiento y sonificación.} El sistema debe mantener una separación funcional entre adquisición, procesamiento de señal y sonificación.

    Esta separación facilita la escalabilidad del sistema y permite modificar un bloque sin reescribir completamente los demás.

\end{description}

\subsubsection{Requisitos no funcionales}

Los requisitos no funcionales recogen las condiciones de calidad, seguridad y operación que debe cumplir el prototipo durante su uso experimental.

\begin{description}[leftmargin=!,labelwidth=2cm]

    \item[RNF-1] \textbf{Funcionamiento en tiempo real.} La cadena de adquisición, procesamiento, sonificación y salida MIDI debe operar de forma continua durante la sesión, sin depender de análisis \textit{offline} posteriores.

    \item[RNF-2] \textbf{Seguridad y uso no clínico.} El sistema debe orientarse a un uso experimental, con alimentación de baja tensión y sin finalidad diagnóstica o clínica.

    \item[RNF-3] \textbf{Robustez frente al ruido.} El diseño debe considerar la baja amplitud del EEG y su sensibilidad a interferencias eléctricas, artefactos fisiológicos y problemas de contacto electrodo--piel.

    \item[RNF-4] \textbf{Modularidad.} El sistema debe separarse en bloques funcionales diferenciados: adquisición, comunicación, procesamiento, sonificación, salida MIDI e interfaz auxiliar.

    \item[RNF-5] \textbf{Escalabilidad.} La arquitectura debe permitir futuras ampliaciones, como nuevos canales EEG, otros mapeos musicales, selección dinámica de canal o nuevas visualizaciones.

    \item[RNF-6] \textbf{Bajo coste y accesibilidad.} El prototipo debe priorizar componentes y herramientas accesibles en un contexto académico y de prototipado.

    \item[RNF-7] \textbf{Replicabilidad y documentación.} El sistema debe documentarse de forma que pueda ser comprendido, reproducido y modificado posteriormente.

    \item[RNF-8] \textbf{Mantenibilidad.} El código debe organizarse de forma clara, reduciendo dependencias innecesarias entre módulos y facilitando ajustes posteriores.

    \item[RNF-9] \textbf{Portabilidad.} El sistema debe poder utilizarse como prototipo compacto y transportable en un entorno experimental controlado.

\end{description}

\subsubsection{Decisiones de diseño adoptadas}

A continuación se resumen las decisiones principales tomadas para implementar la versión desarrollada del sistema.

\begin{description}[leftmargin=!,labelwidth=1.8cm]

    \item[DD-1] \textbf{ADS1299--4PAG.} Se seleccionó el ADS1299--4PAG como front-end de adquisición por estar orientado a biopotenciales y permitir la digitalización multicanal de señales EEG de baja amplitud.

    \item[DD-2] \textbf{Frecuencia de muestreo de 250 Hz.} En las capturas y resultados del TFG se utilizó una frecuencia de \SI{250}{\hertz} por canal, suficiente para el análisis de las bandas EEG empleadas en la sonificación.

    \item[DD-3] \textbf{Montaje ear--EEG/mastoides.} Durante la validación se empleó una configuración próxima a ear--EEG/mastoides, con referencia o BIAS en la muñeca izquierda, por su estabilidad práctica y menor interferencia con el cabello.

    \item[DD-4] \textbf{Procesamiento en Python.} El procesamiento principal se implementó en Python dentro del entorno Linux/Qualcomm de Arduino UNO Q, facilitando el desarrollo y ajuste de los algoritmos.

    \item[DD-5] \textbf{Estimación multitaper.} Se empleó el método multitaper para obtener características espectrales más estables en señales EEG no estacionarias.

    \item[DD-6] \textbf{Ventana de análisis de 4 segundos.} Se utilizó una ventana temporal aproximada de \SI{4}{\second} como compromiso entre estabilidad frecuencial y respuesta temporal.

    \item[DD-7] \textbf{Salida MIDI por D1/RX.} Los eventos MIDI generados se enviaron mediante la línea asociada al pin D1/RX hacia la PCB MIDI del prototipo.

    \item[DD-8] \textbf{Interfaz web auxiliar.} La interfaz web se utilizó como herramienta de monitorización y apoyo durante las pruebas del sistema.

    \item[DD-9] \textbf{Matriz LED opcional.} La matriz LED se consideró un complemento visual opcional para representar la actividad musical generada.

\end{description}

\subsection{Arquitectura general del sistema}

La arquitectura general del sistema se plantea como una cadena modular orientada a transformar una señal EEG real en eventos musicales MIDI enviados por una salida física. El diseño separa de forma explícita las responsabilidades de adquisición, comunicación, procesamiento digital, evaluación de calidad, sonificación, planificación MIDI y monitorización auxiliar. Esta separación permite que cada bloque pueda validarse de forma independiente y evita que la lógica musical dependa directamente de detalles de bajo nivel del convertidor analógico-digital o del firmware.

El flujo principal comienza en los electrodos de superficie, que recogen la actividad bioeléctrica del usuario. Estas señales se introducen en el front-end ADS1299--4PAG, encargado de digitalizar la señal EEG. La lectura del convertidor se realiza desde el microcontrolador de la Arduino UNO Q mediante una comunicación SPI sincronizada con la señal DRDY y el modo de adquisición continua RDATAC. En esta primera etapa, el firmware configura el ADS1299, valida el estado de los frames recibidos, reconstruye las muestras de 24 bits, aplica el escalado correspondiente y entrega la señal en microvoltios.

Una característica importante de la arquitectura es que la comunicación entre el microcontrolador y el entorno Linux no se realiza muestra a muestra de forma aislada, sino mediante bloques de muestras. El contrato principal de transmisión es el evento \texttt{eeg\_block\_uV}, que agrupa bloques de 8 muestras y conserva una estructura de cuatro canales. Aunque la configuración final de validación utiliza CH1 como canal EEG principal para la sonificación, se mantiene \texttt{NUM\_CH=4} en el contrato de datos para preservar la compatibilidad con el backend, las capturas CSV, las herramientas de análisis y la interfaz de monitorización. De esta forma, la decisión experimental de trabajar con un canal principal no rompe la estructura multicanal del sistema.

En la parte Linux de la Arduino UNO Q, el backend Python recibe los bloques enviados por el firmware, valida su continuidad y los introduce en un buffer temporal multicanal. A partir de este buffer se extraen ventanas de señal para el procesamiento digital. El cálculo de características se realiza sobre ventanas temporales, no sobre muestras individuales, lo que permite obtener medidas espectrales más estables. En la implementación validada, el sistema trabaja a \SI{250}{\hertz} y calcula nuevas características cada 64 muestras, con una ventana de análisis de \SI{4}{\second}. Estos valores forman parte de la configuración de diseño empleada en las pruebas, no de una limitación conceptual de la arquitectura.

El procesamiento digital se basa en la extracción de características temporales y espectrales de la señal EEG. Entre ellas se incluyen la amplitud RMS, la potencia absoluta y relativa por bandas EEG, y las frecuencias pico en distintas bandas. La estimación espectral se realiza mediante el bloque \texttt{DSPCore}, que actúa como fuente centralizada para el cálculo de PSD y potencias por banda. Estas características no se utilizan todavía como música, sino como una representación intermedia de la señal sobre la que se puede razonar de forma más estable que sobre la forma de onda cruda.

Antes de alimentar el motor de sonificación, el sistema evalúa la calidad de la ventana analizada mediante un bloque de calidad espectral. Este bloque genera un estado de calidad y un factor de puerta o atenuación. Su función es impedir que ventanas dominadas por artefactos, pérdidas de datos, saturación, ruido excesivo o actividad no fiable modulen la música con la misma intensidad que una ventana válida. Por tanto, el \textit{quality gate} no es únicamente una herramienta de validación posterior, sino un elemento de diseño integrado en la arquitectura de ejecución.

El módulo de sonificación recibe las características EEG ya procesadas y condicionadas por el bloque de calidad. La señal EEG no se convierte directamente en notas musicales ni decide por sí sola la tonalidad del sistema. En la arquitectura final, parámetros musicales como la nota raíz, la nota principal y la escala son controles de configuración musical disponibles para el usuario o para la WebUI. El EEG modula principalmente la actividad, la densidad rítmica, el registro, la estabilidad, la dinámica y la probabilidad de generación de notas. Esta decisión evita una traducción excesivamente literal de EEG a sonido y permite mantener una salida musical más controlada y coherente.

La generación musical se organiza jerárquicamente. Primero se construye un estado musical de alto nivel a partir de los controles derivados del EEG. Después se generan compases, posiciones rítmicas, acordes y finalmente notas individuales. Estas notas se representan como eventos con instante de inicio, duración, altura, velocidad y canal MIDI. Posteriormente, el planificador MIDI transforma dichas notas en eventos temporizados de activación y desactivación, evitando que queden notas colgadas y permitiendo emitir mensajes de seguridad como \textit{panic} cuando sea necesario.

La salida principal del sistema es el MIDI físico. Los eventos MIDI generados en Python se convierten en bytes MIDI y se envían al firmware mediante el contrato \texttt{midi\_bytes}. En el microcontrolador, el handler correspondiente escribe los bytes por la UART asociada a la salida MIDI física del prototipo. Esta ruta permite conectar el sistema a sintetizadores, módulos de sonido o interfaces externas compatibles con MIDI. Por ello, la WebUI no se considera la salida principal del sistema, sino una herramienta auxiliar de supervisión y control.

La interfaz web cumple una función de observación y control ligero. Permite visualizar el estado de adquisición, métricas de recepción, características EEG, diagnóstico de calidad, estado de sonificación, eventos musicales recientes y estado MIDI. También permite modificar parámetros musicales acotados, como la nota raíz, la nota principal y la escala. Sin embargo, no sustituye a la salida MIDI física ni forma parte del flujo esencial de generación de sonido. Su papel es facilitar la supervisión experimental y el ajuste del sistema durante las pruebas.

De forma complementaria, la matriz LED puede actuar como visualización lateral de las notas recientes, reutilizando la misma información que el piano roll de la WebUI. Esta matriz no constituye un flujo musical independiente ni modifica la sonificación; únicamente representa visualmente parte de la actividad musical generada. Por este motivo se considera un elemento opcional y no un requisito del flujo de ejecución principal.

La Figura~\ref{fig:arquitectura_general_sistema} resume esta organización. En ella se diferencia la ruta esencial EEG--MIDI de las herramientas auxiliares de monitorización, visualización y validación.

\begin{figure}[H]
    \centering
    \resizebox{0.98\textwidth}{!}{%
    \begin{tikzpicture}[
        block/.style={draw, rounded corners, align=center, minimum width=2.75cm, minimum height=1.0cm, font=\small},
        aux/.style={draw, rounded corners, dashed, align=center, minimum width=2.75cm, minimum height=1.0cm, font=\small},
        arrow/.style={-{Latex[length=3mm]}, thick},
        auxarrow/.style={-{Latex[length=2.5mm]}, dashed},
        node distance=0.95cm and 0.65cm
    ]
        \node[block] (electrodos) {Electrodos\\EEG};
        \node[block, right=of electrodos] (ads) {ADS1299--4PAG\\adquisición};
        \node[block, right=of ads] (fw) {Firmware MCU\\STM32U585};
        \node[block, right=of fw] (bridge) {\texttt{eeg\_block\_uV}\\bloques EEG};

        \node[block, below=1.15cm of bridge] (backend) {Backend\\Python/Linux};
        \node[block, left=of backend] (dsp) {DSPCore\\features EEG};
        \node[block, left=of dsp] (quality) {Quality gate\\calidad};
        \node[block, left=of quality] (sonif) {Motor de\\sonificación};
        \node[block, below=1.15cm of sonif] (midi) {Scheduler y\\transporte MIDI};
        \node[block, below=1.15cm of midi] (midiout) {MIDI OUT\\físico};

        \node[aux, below=1.15cm of dsp] (web) {WebUI\\monitorización};
        \node[aux, below=1.15cm of quality] (led) {Piano roll / LED\\opcional};
        \node[aux, right=of backend] (tools) {Capturas y\\benchmarks};

        \draw[arrow] (electrodos) -- (ads);
        \draw[arrow] (ads) -- node[above, font=\scriptsize]{SPI} (fw);
        \draw[arrow] (fw) -- (bridge);
        \draw[arrow] (bridge) -- (backend);
        \draw[arrow] (backend) -- (dsp);
        \draw[arrow] (dsp) -- (quality);
        \draw[arrow] (quality) -- (sonif);
        \draw[arrow] (sonif) -- (midi);
        \draw[arrow] (midi) -- node[right, font=\scriptsize]{\texttt{UART}} (midiout);

        \draw[auxarrow] (backend) -- (web);
        \draw[auxarrow] (sonif) -- (led);
        \draw[auxarrow] (backend) -- (tools);
    \end{tikzpicture}}
    \caption{Arquitectura general del sistema EEG--MIDI. La ruta principal transforma EEG real en salida MIDI física; la WebUI, la matriz LED, las capturas y los benchmarks actúan como elementos auxiliares de monitorización, visualización y validación.}
    \label{fig:arquitectura_general_sistema}
\end{figure}

En conjunto, la arquitectura permite distinguir tres niveles. El primer nivel corresponde al tiempo real embebido, donde el microcontrolador adquiere, filtra y empaqueta muestras EEG. El segundo nivel corresponde al procesamiento y generación musical en Python, donde se calculan características, se evalúa la calidad y se generan eventos MIDI. El tercer nivel corresponde a herramientas auxiliares de observación, captura, benchmark y análisis \textit{offline}, que aportan trazabilidad y validación, pero no forman parte de la cadena esencial EEG--MIDI. Esta separación resulta clave para documentar el sistema de forma rigurosa y para justificar posteriormente los apartados de implementación y validación.

\subsection{Bloque de adquisición EEG}

El bloque de adquisición EEG constituye la primera etapa funcional del sistema. Su responsabilidad es captar la actividad bioeléctrica del usuario mediante electrodos de superficie, digitalizarla con un \textit{front-end} adecuado para biopotenciales y entregar al firmware una secuencia de muestras temporales coherente, escalada y sincronizada. Esta etapa condiciona al resto de la arquitectura, ya que cualquier error de sincronismo, escala, referencia corporal o configuración de entrada se propagaría posteriormente hacia el procesamiento espectral, el cálculo de calidad y la sonificación MIDI.

La elección del ADS1299--4PAG responde a los requisitos propios de una adquisición EEG. Se trata de un convertidor analógico-digital de 24 bits orientado a biopotenciales, con entradas diferenciales, amplificador de ganancia programable, referencia interna, interfaz SPI, circuitería de BIAS y funciones de diagnóstico como \textit{lead-off} y señal de test interna, según la descripción funcional del fabricante \cite{TexasInstrumentsADS1299}. Estas características permiten integrar en un único componente gran parte del \textit{front-end} necesario para medir señales EEG de baja amplitud, reduciendo la necesidad de etapas externas de amplificación antes de la digitalización.

Desde el punto de vista de diseño, el ADS1299 se utiliza como frontera entre el dominio bioeléctrico y el dominio digital. En el lado analógico se encuentran los electrodos, el par diferencial de medida y el electrodo de BIAS/RLD. En el lado digital se encuentran la interfaz SPI, la señal DRDY, el modo de lectura continua RDATAC y el firmware del microcontrolador. Esta separación permite tratar de forma independiente las decisiones de montaje físico y las decisiones de comunicación digital.

En la configuración final del prototipo se utiliza CH1 como canal EEG principal. El montaje empleado se plantea como una configuración tipo ear--EEG/mastoides, con el par diferencial conectado al canal CH1 y con BIAS/RLD derivado de CH1P y CH1N. Esta configuración busca mejorar la estabilidad de modo común y reducir la influencia de entradas no utilizadas. Las entradas CH2, CH3 y CH4 no se interpretan como EEG activo en las capturas finales; sin embargo, se conservan dentro del contrato de datos para mantener una estructura multicanal estable entre firmware, backend Python, capturas CSV, herramientas de análisis e interfaz de monitorización.

La conservación de cuatro canales en el contrato no debe confundirse con una sonificación multicanal completa. El diseño mantiene la capacidad estructural multicanal del sistema, pero la versión final de validación utiliza CH1 como canal principal para el análisis espectral y la generación musical. Esta decisión permite obtener una señal más controlada durante la validación sin renunciar a una arquitectura ampliable a otros canales en futuras iteraciones.

La configuración de BIAS/RLD se justifica por la sensibilidad de las señales EEG al ruido de modo común y a la referencia corporal. El ADS1299 permite derivar la señal de BIAS a partir de canales seleccionados y realimentarla hacia el usuario mediante la red analógica prevista en la placa. En el sistema final, la derivación se limita al par usado por CH1, evitando incluir canales no utilizados o flotantes. Además, durante la captura final se desactiva la detección \textit{lead-off} para evitar que la corriente de diagnóstico asociada a esa función pueda introducir componentes no deseadas en la señal usada para sonificación.

La adquisición se realiza a una frecuencia de \SI{250}{\hertz} por canal. Esta frecuencia es adecuada para el objetivo del prototipo, ya que cubre las bandas EEG empleadas posteriormente en el procesamiento espectral y mantiene una carga razonable para el firmware, el transporte de datos y el backend Python. La frecuencia de muestreo es un parámetro de coherencia global: debe coincidir entre el firmware, el parser de datos, el buffer temporal, las ventanas de análisis y las herramientas de validación.

La lectura digital del ADS1299 se organiza mediante la señal DRDY y el modo RDATAC. DRDY actúa como indicación hardware de que existe una nueva conversión disponible. El firmware no interroga el convertidor de forma arbitraria, sino que espera esa señal para leer el frame correspondiente por SPI. El uso de RDATAC permite mantener una lectura continua, adecuada para un sistema de funcionamiento en vivo, en el que la señal EEG debe fluir de forma sostenida hacia el backend.

Cada frame leído contiene un campo de estado y las muestras de los canales. El firmware comprueba el prefijo del \textit{status} para detectar posibles frames incoherentes o desalineados antes de aceptar la muestra. A continuación reconstruye los valores de 24 bits con signo, respetando el formato en complemento a dos, y aplica el factor de conversión correspondiente para pasar de cuentas digitales a voltios. Esta conversión es esencial para que la señal conserve una escala física interpretable y para que las herramientas posteriores puedan expresar las muestras en microvoltios.

Tras la reconstrucción y el escalado inicial, el firmware aplica una cadena de filtrado básica en el microcontrolador. Esta cadena incluye un bloqueo de continua o paso alto, un filtro de rechazo de red a \SI{50}{\hertz} y un paso bajo orientado a limitar el contenido de alta frecuencia antes de enviar la señal al entorno Python. La señal transmitida al backend ya se encuentra, por tanto, filtrada de forma básica y expresada en microvoltios.

La Figura~\ref{fig:bloque_adquisicion_eeg} resume el recorrido de la señal dentro del bloque de adquisición, desde los electrodos hasta la obtención de muestras EEG en microvoltios listas para ser empaquetadas por el bloque de comunicación.

\begin{figure}[H]
    \centering
    \begin{tikzpicture}[
        block/.style={draw, rounded corners, align=center, text width=2.65cm, minimum height=0.95cm, font=\footnotesize},
        signal/.style={draw, rounded corners, align=center, text width=2.65cm, minimum height=0.95cm, font=\footnotesize, fill=gray!8},
        arrow/.style={-{Latex[length=3mm]}, thick},
        note/.style={align=center, font=\scriptsize}
    ]
        \node[block]  (electrodos) at (0,0) {Electrodos\\de superficie};
        \node[block]  (montaje)   at (3.4,0) {CH1P / CH1N\\BIAS/RLD};
        \node[block]  (ads)       at (6.8,0) {ADS1299--4PAG\\AFE 24 bits};
        \node[block]  (digital)   at (10.2,0) {SPI\\DRDY + RDATAC};

        \node[signal] (frame)     at (10.2,-2.75) {status 24 bits\\CH1--CH4 signed 24 bits};
        \node[signal] (escala)    at (6.8,-2.75) {validación de status\\escalado \texttt{LSB\_V}};
        \node[signal] (filtros)   at (3.4,-2.75) {filtros MCU\\HP + notch 50 Hz + LP};
        \node[signal] (uv)        at (0,-2.75) {muestras EEG\\en \si{\micro\volt}};

        \draw[arrow] (electrodos) -- (montaje);
        \draw[arrow] (montaje) -- (ads);
        \draw[arrow] (ads) -- (digital);
        \draw[arrow] (digital) -- (frame);
        \draw[arrow] (frame) -- (escala);
        \draw[arrow] (escala) -- (filtros);
        \draw[arrow] (filtros) -- (uv);

        \node[note] at (3.4,-1.1) {CH1 es el canal EEG principal\\en la configuración final};
        \node[note] at (6.8,-4.12) {CH2--CH4 se conservan\\por contrato de datos};
    \end{tikzpicture}
    \caption{Bloque de adquisición EEG basado en el ADS1299--4PAG. La etapa transforma la señal bioeléctrica en muestras digitales escaladas en microvoltios; el empaquetado en bloques \texttt{eeg\_block\_uV} pertenece al bloque de comunicación posterior.}
    \label{fig:bloque_adquisicion_eeg}
\end{figure}

El bloque de adquisición finaliza cuando las muestras EEG quedan disponibles en el firmware como valores temporales en microvoltios, asociados a su \textit{status} y a su índice de muestra. A partir de ese punto comienza el bloque de comunicación microcontrolador--Linux, encargado de agrupar las muestras en bloques y transmitirlas mediante el evento \texttt{eeg\_block\_uV}. Por tanto, el bloque de adquisición no calcula potencias por banda, métricas de calidad ni parámetros musicales; su función es obtener una señal digitalizada, sincronizada, validada y escalada correctamente.

En resumen, el bloque de adquisición EEG se define por la siguiente cadena: electrodos de superficie, par diferencial CH1, BIAS/RLD, ADS1299--4PAG, lectura SPI sincronizada por DRDY, modo RDATAC, reconstrucción de datos de 24 bits, validación de \textit{status}, conversión a unidades físicas, filtrado básico en MCU y salida en microvoltios. Esta delimitación permite justificar de forma clara qué decisiones pertenecen al diseño de adquisición y cuáles corresponden a los bloques posteriores de comunicación, procesamiento digital, calidad espectral y sonificación.

\subsection{Bloque de comunicación microcontrolador--Linux}

El bloque de comunicación microcontrolador--Linux conecta la adquisición embebida con el procesamiento en Python. Su función es transportar de forma continua las muestras EEG ya digitalizadas, filtradas y escaladas en microvoltios desde el firmware del STM32U585 hasta el backend que se ejecuta en el entorno Linux de la Arduino UNO Q. Este bloque no modifica el contenido fisiológico de la señal ni calcula características espectrales; su responsabilidad principal es preservar la continuidad temporal, la estructura multicanal y la integridad del contrato de datos.

La comunicación no se realiza muestra a muestra, sino mediante bloques de muestras. Esta decisión reduce el número de llamadas entre el microcontrolador y el entorno Linux, disminuye la sobrecarga asociada al puente de comunicación y mantiene un flujo regular hacia Python. En la configuración final, cada bloque agrupa hasta 8 muestras consecutivas. Dado que la adquisición se realiza a \SI{250}{\hertz}, esta agrupación genera una frecuencia esperada de aproximadamente 31,25 bloques por segundo, suficientemente fina para mantener la continuidad temporal del sistema y, al mismo tiempo, más eficiente que una notificación por muestra individual.

El contrato principal entre firmware y Python se denomina \texttt{eeg\_block\_uV}. Este evento transporta un índice de bloque, el índice de la primera muestra, el número de muestras válidas y una secuencia de valores asociados a cada muestra. Cada muestra conserva el \textit{status} del ADS1299 y los valores de los cuatro canales en microvoltios. Aunque la sonificación final utiliza CH1 como canal EEG principal, el \textit{payload} mantiene una estructura de cuatro canales para preservar la compatibilidad con el backend, las capturas CSV, las herramientas de análisis y la interfaz de monitorización.

La inclusión de \texttt{block\_idx} y \texttt{first\_sample\_idx} permite reconstruir la continuidad de la señal en el lado Python. El índice de bloque permite detectar pérdidas o saltos entre bloques sucesivos. El índice de primera muestra permite verificar la progresión temporal dentro de la secuencia de adquisición. De este modo, el sistema puede distinguir entre una señal correctamente recibida y una señal con discontinuidades, pérdidas o errores de transporte.

El \textit{status} del ADS1299 también se mantiene dentro del contrato. Esta decisión permite que Python no reciba únicamente valores de señal, sino también información mínima de sincronización del convertidor. El receptor puede comprobar que el prefijo de \textit{status} coincide con el formato esperado antes de aceptar el bloque como válido. Esta comprobación actúa como una barrera temprana frente a frames desalineados, \textit{payloads} corruptos o errores de comunicación.

En el firmware, las muestras se acumulan en una cola de transmisión antes de publicarse hacia Linux. Esta cola desacopla parcialmente el ritmo de adquisición, marcado por DRDY, del ritmo de publicación mediante Bridge. El firmware debe seguir atendiendo la adquisición a \SI{250}{\hertz}, por lo que la comunicación no debe bloquear ni introducir trabajo pesado dentro de la ruta crítica de captura. Por este motivo, el bloque de comunicación se diseña como una etapa de empaquetado y publicación, no como una etapa de análisis.

En el lado Python, el evento \texttt{eeg\_block\_uV} es recibido por un componente específico de recepción. Este receptor valida el bloque antes de entregarlo al procesamiento digital. Las comprobaciones principales incluyen el número de muestras, la longitud del \textit{payload}, la continuidad de los índices y el prefijo de \textit{status} ADS1299. Si el bloque cumple el contrato, se encola para ser drenado posteriormente por el backend. Esta separación evita que el \textit{callback} de recepción realice cálculos costosos y permite mantenerlo como una función ligera.

El backend Python drena periódicamente la cola de bloques recibidos y entrega las muestras al procesador EEG. En esta transición, los valores en microvoltios pasan al buffer temporal multicanal utilizado por el bloque DSP. Si existe una captura activa, los mismos bloques también pueden enviarse al gestor de capturas, pero esta ruta es lateral y no forma parte del flujo mínimo necesario para producir MIDI en vivo.

La Figura~\ref{fig:comunicacion_mcu_linux} resume la secuencia principal de comunicación entre el firmware y el backend Python. La figura muestra la ruta funcional de los bloques EEG y separa explícitamente la recepción en tiempo real de la ruta lateral de captura.

\begin{figure}[H]
    \centering
    \begin{tikzpicture}[
        commblock/.style={draw, rounded corners, align=center, text width=4.3cm, minimum height=0.9cm, font=\footnotesize},
        payloadblock/.style={draw, rounded corners, align=center, text width=5.65cm, minimum height=1.1cm, font=\footnotesize, fill=gray!8},
        auxblock/.style={draw, rounded corners, dashed, align=center, text width=3.3cm, minimum height=0.85cm, font=\footnotesize, fill=gray!5},
        arrow/.style={-{Latex[length=3mm]}, thick},
        auxarrow/.style={-{Latex[length=2.5mm]}, dashed, thick},
        note/.style={align=center, font=\scriptsize}
    ]
        \node[commblock]    (mcu)     at (0,0) {Firmware MCU\\STM32U585};
        \node[commblock]    (ring)    at (0,-2.35) {\texttt{TxBlockRing}\\agrupa 8 muestras};
        \node[payloadblock] (notify)  at (0,-5.25) {\texttt{Bridge.notify(\"eeg\_block\_uV\")}\\\texttt{block\_idx}, \texttt{first\_sample\_idx}, \texttt{sample\_count}\\\texttt{status + ch1..ch4} en \si{\micro\volt}};
        \node[commblock]    (rx)      at (0,-7.65) {\texttt{EEGReceiver}\\validación de contrato};
        \node[commblock]    (queue)   at (0,-9.35) {cola RX\\bloques válidos};
        \node[commblock]    (backend) at (0,-11.35) {\texttt{BackendService}\\drenado periódico};
        \node[commblock]    (proc)    at (0,-13.7) {\texttt{EEGSignalProcessor}\\buffer multicanal};
        \node[auxblock]     (capture) at (5.0,-9.35) {\texttt{CaptureManager}\\ruta lateral};

        \draw[arrow] (mcu) -- (ring);
        \draw[arrow] (ring) -- (notify);
        \draw[arrow] (notify) -- (rx);
        \draw[arrow] (rx) -- (queue);
        \draw[arrow] (queue) -- (backend);
        \draw[arrow] (backend) -- (proc);
        \draw[auxarrow] (queue) -- (capture);

        \node[note] at (-4.2,-2.35) {Adquisición a \SI{250}{\hertz}\\31,25 bloques/s esperados};
        \node[note] at (4.2,-7.65) {Comprueba longitud,\\índices y status ADS1299};
        \node[note] at (4.2,-11.7) {Salida del bloque:\\muestras listas para DSP};
    \end{tikzpicture}
    \caption{Secuencia simplificada de comunicación microcontrolador--Linux. El firmware empaqueta muestras EEG en bloques, el Bridge transporta el evento \texttt{eeg\_block\_uV}, Python valida el contrato y el backend drena los bloques hacia el procesador EEG. La captura se representa como ruta lateral de validación.}
    \label{fig:comunicacion_mcu_linux}
\end{figure}

El límite funcional de este bloque se sitúa antes del cálculo de características. La comunicación microcontrolador--Linux no estima bandas EEG, no calcula calidad espectral y no genera eventos musicales. Su salida esperada es una secuencia de bloques EEG válidos, continuos y ordenados, listos para ser incorporados al buffer de procesamiento. A partir de ese punto comienza el bloque de procesamiento digital.

En conjunto, este diseño permite separar tres responsabilidades: el firmware adquiere y empaqueta la señal, el Bridge transporta bloques de datos entre MCU y Linux, y Python valida y encola esos bloques antes de introducirlos en el procesador EEG. Esta separación facilita la detección de errores, permite medir pérdidas o discontinuidades y mantiene estable el contrato entre adquisición, procesamiento, capturas y monitorización.

\subsection{Bloque de procesamiento digital}

El bloque de procesamiento digital transforma los bloques EEG recibidos desde el firmware en características temporales y espectrales utilizables por la sonificación. Su entrada son bloques de muestras EEG ya validados por el receptor Python, expresados en microvoltios y organizados según el contrato multicanal del sistema. Su salida no son todavía eventos musicales, sino una representación intermedia de la señal formada por métricas de amplitud, potencia por bandas, frecuencias pico y estado de calidad.

A diferencia del bloque de adquisición, el procesamiento digital no accede directamente al ADS1299 ni a la interfaz SPI. Tampoco modifica el contrato de comunicación entre el microcontrolador y Linux. Su función comienza cuando el backend Python drena la cola de bloques válidos y entrega las muestras al procesador EEG. En ese punto, \texttt{EEGSignalProcessor} convierte los valores de microvoltios a voltios y los inserta en un buffer circular multicanal. Esta conversión de unidades es importante porque el firmware y las capturas trabajan en microvoltios, mientras que los cálculos internos de PSD y RMS se realizan sobre valores en voltios.

El uso de un buffer circular permite mantener una historia temporal reciente de la señal sin crecer indefinidamente en memoria. En la configuración final, el backend conserva varios segundos de señal multicanal, suficientes para extraer ventanas de análisis y actualizar las características de forma periódica. Aunque el contrato mantiene cuatro canales, la versión final del sistema utiliza CH1 como canal principal para el análisis espectral en vivo y para la sonificación. Esta decisión simplifica la validación experimental sin eliminar la estructura multicanal del sistema.

El cálculo de características no se realiza muestra a muestra. El sistema espera a disponer de una ventana temporal suficiente y calcula nuevas características cada cierto número de muestras. En la configuración final, la ventana de análisis es de \SI{4}{\second} y el salto entre actualizaciones es de 64 muestras. Con una frecuencia de muestreo de \SI{250}{\hertz}, esto equivale a una actualización aproximada cada \SI{256}{\milli\second}. Esta elección busca un compromiso entre estabilidad espectral y respuesta temporal: una ventana demasiado corta produciría estimaciones frecuenciales poco estables, mientras que una ventana excesivamente larga retrasaría la respuesta musical del sistema.

El núcleo del análisis espectral se concentra en \texttt{DSPCore}. Este bloque calcula la densidad espectral de potencia de la ventana EEG y obtiene métricas asociadas a las bandas delta, theta, alfa, beta y gamma. En la ruta en vivo se emplea estimación multitaper, ya que ofrece una estimación espectral más estable para señales no estacionarias que una única ventana clásica. A partir de la PSD se calculan potencias absolutas, potencias relativas y frecuencias pico. Además, se calcula la amplitud RMS de la señal, que aporta una medida temporal complementaria a la información frecuencial.

Las potencias relativas por banda son especialmente relevantes para la sonificación, porque permiten describir la distribución espectral de la señal con menor dependencia de la amplitud absoluta. Las potencias absolutas y el RMS, por otro lado, resultan útiles para diagnóstico, trazabilidad y detección de cambios de energía en la señal. Las frecuencias pico permiten identificar la componente dominante global o dentro de bandas concretas, aunque deben interpretarse con cautela cuando existe ruido, artefactos musculares o mala calidad de contacto.

Después del cálculo de características, el sistema evalúa la calidad de la ventana analizada. Esta etapa calcula métricas de diagnóstico como amplitud RMS en microvoltios, rango pico a pico, presencia de componente de red, saturación, discontinuidades o comportamiento anómalo de la forma de onda. Estas métricas se combinan posteriormente en un bloque de calidad espectral que produce un estado de calidad y un factor de puerta. El objetivo es evitar que ventanas dominadas por ruido, artefactos o errores de transporte influyan en la sonificación con la misma intensidad que una ventana limpia.

El \textit{quality gate} actúa, por tanto, como frontera entre procesamiento EEG y sonificación. Si la ventana se considera válida, las características pueden modular la generación musical con normalidad. Si la calidad es dudosa o mala, el sistema puede atenuar o bloquear la influencia de la señal sobre los controles musicales. Esta decisión es importante porque en una aplicación EEG real no todas las ventanas capturadas contienen información fisiológica útil; algunas pueden estar dominadas por movimiento, contacto deficiente, actividad muscular o interferencias.

La Figura~\ref{fig:bloque_procesamiento_digital} resume el flujo interno del bloque de procesamiento digital. La figura diferencia entre la ruta principal de cálculo de características y la evaluación de calidad que condiciona su uso posterior en la sonificación.

\begin{figure}[H]
    \centering
    \begin{tikzpicture}[
        procblock/.style={draw, rounded corners, align=center, text width=3.9cm, minimum height=0.95cm, font=\footnotesize},
        featureblock/.style={draw, rounded corners, align=center, text width=4.2cm, minimum height=0.95cm, font=\footnotesize, fill=gray!8},
        qualityblock/.style={draw, rounded corners, dashed, align=center, text width=3.9cm, minimum height=0.95cm, font=\footnotesize, fill=gray!5},
        arrow/.style={-{Latex[length=3mm]}, thick},
        auxarrow/.style={-{Latex[length=2.5mm]}, dashed, thick},
        note/.style={align=center, font=\scriptsize}
    ]
        \node[procblock]    (blocks)   at (0,0)        {bloques EEG válidos\\en \si{\micro\volt}};
        \node[procblock]    (units)    at (0,-1.7)     {conversión de unidades\\\si{\micro\volt} $\rightarrow$ V};
        \node[procblock]    (buffer)   at (0,-3.4)     {ring buffer\\multicanal};
        \node[procblock]    (window)   at (0,-5.1)     {ventana CH1\\\SI{4}{\second}};
        \node[featureblock] (dspcore)  at (0,-6.95)    {\texttt{DSPCore}\\PSD multitaper};
        \node[featureblock] (features) at (0,-8.8)     {RMS, bandpowers\\relativos/absolutos y picos};
        \node[qualityblock] (quality)  at (5.8,-8.8)   {diagnóstico\\y \textit{quality gate}};
        \node[featureblock] (out)      at (0,-10.65)   {features EEG\\para sonificación};

        \draw[arrow] (blocks) -- (units);
        \draw[arrow] (units) -- (buffer);
        \draw[arrow] (buffer) -- (window);
        \draw[arrow] (window) -- (dspcore);
        \draw[arrow] (dspcore) -- (features);
        \draw[arrow] (features) -- (out);

        \draw[auxarrow] (features) -- (quality);
        \draw[auxarrow] (quality) |- (out);

        \node[note] at (-3.7,-5.1)  {Actualización cada\\64 muestras};
        \node[note] at (-3.7,-6.95) {Con \SI{250}{\hertz}:\\aprox. \SI{256}{\milli\second}};
    \end{tikzpicture}
    \caption{Flujo del bloque de procesamiento digital. Python convierte los bloques EEG a voltios, mantiene un buffer circular, extrae una ventana de CH1, calcula características espectrales mediante \texttt{DSPCore} y evalúa la calidad de la ventana antes de entregar las features al bloque de sonificación.}
    \label{fig:bloque_procesamiento_digital}
\end{figure}

El bloque de procesamiento digital finaliza generando una estructura de características y diagnóstico que será consumida por el adaptador de sonificación. Esta salida mantiene separadas dos capas: por un lado, las características EEG calculadas a partir de la señal; por otro, la evaluación de calidad que determina cuánto debe confiar el sistema en esa ventana. La conversión de dichas características en parámetros musicales pertenece al bloque posterior de sonificación EEG--MIDI.

En resumen, el procesamiento digital se define como la cadena formada por recepción de bloques válidos, conversión de microvoltios a voltios, almacenamiento en buffer circular multicanal, extracción de una ventana temporal de CH1, estimación espectral multitaper, cálculo de RMS, potencias por banda y picos, diagnóstico de calidad y generación del factor de puerta. Esta organización permite que la sonificación no dependa directamente de la señal cruda, sino de una representación EEG más estable, interpretable y controlada.

\subsection{Bloque de sonificación EEG--MIDI}

El bloque de sonificación EEG--MIDI transforma las características calculadas a partir de la señal EEG en eventos musicales enviados posteriormente por la salida MIDI física del sistema. Su entrada no es la señal EEG cruda, sino una representación ya procesada formada por potencias por bandas, RMS, frecuencias pico y estado de calidad. Esta decisión separa claramente el análisis fisiológico de la generación musical y evita que la música dependa directamente de muestras individuales o de fluctuaciones instantáneas de la señal.

La sonificación se diseña como una etapa jerárquica. En primer lugar, las características EEG y el resultado del bloque de calidad se convierten en controles musicales normalizados. Estos controles describen aspectos musicales como nivel de actividad, calma, tensión, densidad rítmica, registro, estabilidad armónica, intensidad y probabilidad de generación de notas. Aunque estos nombres se utilizan internamente en el sistema, deben interpretarse como variables musicales derivadas de la señal EEG, no como medidas clínicas ni diagnósticas.

El uso del \textit{quality gate} es fundamental en esta etapa. La sonificación no debe responder con la misma intensidad a una ventana EEG válida que a una ventana dominada por ruido, artefactos musculares, saturación, mala conexión de electrodos o discontinuidades de transporte. Por ello, el estado de calidad puede atenuar o bloquear la influencia de las características EEG sobre los controles musicales. Esta barrera evita que artefactos no fisiológicos produzcan cambios musicales bruscos o engañosos.

Una vez obtenidos los controles musicales, el sistema no genera notas de forma aislada. La generación se organiza en niveles. Primero se construye un estado musical de alto nivel, que resume la intención musical de la ventana EEG actual. Después se genera una estructura de compás, se seleccionan acordes o contextos armónicos y se distribuyen posiciones rítmicas. Finalmente se crean eventos de nota individuales, cada uno con altura, inicio, duración, velocidad y canal MIDI.

Esta organización permite que la salida musical tenga mayor coherencia temporal que un mapeo directo banda--nota. En lugar de disparar notas de forma inmediata por cada variación espectral, el sistema acumula la información EEG en controles que modifican gradualmente la densidad, el registro, la dinámica y la estabilidad musical. Así se reduce la probabilidad de obtener una secuencia caótica o excesivamente sensible al ruido de la señal.

La arquitectura también separa los parámetros musicales de configuración de los parámetros derivados del EEG. La nota raíz, la nota principal y la escala pueden modificarse desde la interfaz auxiliar, pero estos cambios no alteran la adquisición, el DSP ni el transporte MIDI. Su función es definir el marco musical sobre el que actúan las características EEG. De esta forma, el usuario puede ajustar el carácter armónico del sistema sin romper el flujo principal EEG--MIDI.

Una vez generadas las notas, el planificador MIDI transforma cada nota en eventos temporizados. Una nota musical no se representa únicamente como una activación, sino como una pareja de eventos \textit{Note On} y \textit{Note Off}. El planificador mantiene el control de las notas activas, ordena los eventos pendientes y permite liberar notas de forma segura. Esta etapa es necesaria para evitar notas colgadas y para mantener una salida MIDI estable durante la ejecución continua.

El transporte MIDI convierte los eventos planificados en bytes MIDI y los envía al firmware mediante el contrato \texttt{midi\_bytes}. En el microcontrolador, estos bytes se escriben por la UART asociada a la salida MIDI física del prototipo. Por tanto, la salida principal del bloque no es una gráfica ni una simulación en pantalla, sino una secuencia real de mensajes MIDI enviada hacia el conector físico MIDI OUT.

La WebUI y el piano roll cumplen una función auxiliar. Permiten observar las notas recientes, el estado de la sonificación, los parámetros musicales activos y el estado del transporte MIDI. Sin embargo, no sustituyen al flujo principal. La función central del sistema sigue siendo transformar actividad EEG real en eventos musicales enviados por MIDI físico.

La matriz LED, cuando está habilitada, también se considera una visualización lateral. Reutiliza información de notas recientes para representar actividad musical, pero no interviene en la generación de notas ni en la planificación MIDI. Por tanto, ni la WebUI ni la matriz LED forman parte de la ruta mínima necesaria para que el sistema produzca sonido mediante MIDI.

La Figura~\ref{fig:bloque_sonificacion_midi} resume la organización interna del bloque de sonificación y su conexión con la salida MIDI física. La figura muestra el flujo principal de generación musical y sitúa la WebUI como ruta auxiliar de configuración y monitorización.

\begin{figure}[H]
    \centering
    \begin{tikzpicture}[
        sonblock/.style={draw, rounded corners, align=center, text width=4.3cm, minimum height=0.95cm, font=\footnotesize},
        musicblock/.style={draw, rounded corners, align=center, text width=4.3cm, minimum height=0.95cm, font=\footnotesize, fill=gray!8},
        auxblock/.style={draw, rounded corners, dashed, align=center, text width=3.7cm, minimum height=0.9cm, font=\footnotesize, fill=gray!5},
        arrow/.style={-{Latex[length=3mm]}, thick},
        auxarrow/.style={-{Latex[length=2.5mm]}, dashed, thick},
        note/.style={align=center, font=\scriptsize}
    ]
        \node[sonblock]   (features)  at (0,0)      {features EEG\\+ \textit{quality gate}};
        \node[sonblock]   (adapter)   at (0,-1.55)  {\texttt{SonificationFeatureAdapter}\\controles musicales};
        \node[musicblock] (segment)   at (0,-3.1)   {\texttt{MusicSegmentBuilder}\\estado musical};
        \node[musicblock] (bar)       at (0,-4.65)  {\texttt{BarGenerator}\\compás y contexto armónico};
        \node[musicblock] (notes)     at (0,-6.2)   {\texttt{NoteGenerator}\\eventos de nota};
        \node[musicblock] (scheduler) at (0,-7.75)  {\texttt{MidiScheduler}\\\textit{Note On/Off}};
        \node[sonblock]   (transport) at (0,-9.3)   {\texttt{MidiByteTransport}\\\texttt{Bridge.call("midi\_bytes")}};
        \node[sonblock]   (midiout)   at (0,-10.85) {MIDI OUT\\físico};

        \node[auxblock] (web) at (5.4,-3.1) {WebUI auxiliar\\root, main note y escala};
        \node[auxblock] (led) at (5.4,-6.2) {Piano roll / LED\\visualización lateral};

        \draw[arrow] (features) -- (adapter);
        \draw[arrow] (adapter) -- (segment);
        \draw[arrow] (segment) -- (bar);
        \draw[arrow] (bar) -- (notes);
        \draw[arrow] (notes) -- (scheduler);
        \draw[arrow] (scheduler) -- (transport);
        \draw[arrow] (transport) -- (midiout);

        \draw[auxarrow] (web) -- (segment);
        \draw[auxarrow] (notes) -- (led);

        \node[note] at (-4.4,-1.55) {Los controles no son\\variables clínicas};
        \node[note] at (-4.4,-7.75) {Evita notas colgadas\\mediante \textit{Note Off}/panic};
    \end{tikzpicture}
    \caption{Bloque de sonificación EEG--MIDI. Las características EEG y el \textit{quality gate} se transforman en controles musicales, que alimentan una generación jerárquica de segmentos, compases y notas. El planificador MIDI convierte las notas en eventos temporizados y el transporte envía bytes al firmware para producir la salida MIDI física.}
    \label{fig:bloque_sonificacion_midi}
\end{figure}

En resumen, el bloque de sonificación EEG--MIDI se define como la cadena formada por características EEG, evaluación de calidad, adaptación a controles musicales, construcción de segmento, generación de compás, generación de notas, planificación temporal MIDI y transporte hacia el firmware. Esta separación permite mantener la sonificación controlada, evitar una traducción excesivamente literal de la señal EEG y asegurar que la salida final del sistema sea MIDI físico operativo.

\subsection{Bloque de visualización e interacción}

El bloque de visualización e interacción proporciona una interfaz auxiliar para observar el estado del sistema y modificar parámetros musicales de alto nivel durante la ejecución. No forma parte de la ruta mínima necesaria para producir sonido, ya que la salida principal del prototipo es el puerto MIDI físico. Su función es facilitar la supervisión experimental, la depuración y el ajuste musical sin intervenir directamente en la adquisición, el procesamiento digital ni la escritura de bytes MIDI en la UART.

La interfaz utilizada en la versión final no se basa en Streamlit, sino en una WebUI ligera servida desde el entorno Python de la Arduino UNO Q. El backend genera periódicamente un \textit{snapshot} de estado que resume la adquisición, la recepción de bloques, el buffer EEG, las características espectrales, el diagnóstico de calidad, la sonificación, el estado MIDI, las notas recientes y la configuración auxiliar. Este \textit{snapshot} actúa como contrato de visualización entre el backend y la interfaz.

Desde el punto de vista de diseño, la WebUI se comporta como un observador del sistema. El navegador recibe el estado publicado por el servidor web y actualiza los paneles de visualización. La interfaz no calcula PSD, no estima potencias por banda, no accede al ADS1299, no modifica los filtros del firmware y no genera notas musicales por sí misma. Los cálculos críticos siguen perteneciendo a los bloques de adquisición, comunicación, procesamiento, calidad y sonificación.

La información mostrada se organiza en paneles funcionales. Un primer conjunto resume el estado general de la aplicación, la disponibilidad de ventana de análisis, las tasas de recepción, los índices de muestra, posibles pérdidas de bloques y el estado de los canales. Otro conjunto representa las características EEG de CH1, incluyendo RMS, frecuencias pico y potencias por bandas. También se muestran métricas de diagnóstico y calidad, como amplitud, rango pico a pico, componente de red, avisos y estado del \textit{quality gate}.

La WebUI incorpora además paneles específicos de sonificación y MIDI. Estos paneles permiten observar los controles musicales derivados de las características EEG, el estado del planificador, el transporte MIDI, el estado del handler en el microcontrolador y las notas recientes generadas. El piano roll representa visualmente dichas notas en el tiempo y actúa como evidencia rápida de que el motor de sonificación está produciendo eventos musicales.

La interacción del usuario se mantiene deliberadamente acotada. La interfaz permite modificar parámetros musicales como la nota raíz, la nota principal y la escala. Estos controles cambian el marco musical sobre el que opera la sonificación, pero no alteran el firmware, el ADS1299, los filtros, la frecuencia de muestreo, el contrato \texttt{eeg\_block\_uV} ni el contrato \texttt{midi\_bytes}. Esta separación reduce el riesgo de que una acción de interfaz deje el sistema en un estado de adquisición no comparable o difícil de validar.

Otra interacción importante es el botón de \textit{panic} MIDI. Su finalidad es enviar una acción de seguridad cuando existe riesgo de notas activas no liberadas o cuando se desea limpiar el estado musical durante las pruebas. Aunque esta acción es auxiliar, resulta relevante para la operabilidad del sistema porque permite recuperar una salida MIDI estable sin reiniciar toda la aplicación.

Algunas rutas de la WebUI tienen carácter diagnóstico, como las pruebas MIDI de nota, secuencia o bucle. Estas funciones son útiles para comprobar el transporte MIDI de forma aislada, pero no pertenecen al flujo principal EEG--MIDI. En la memoria se consideran herramientas de apoyo y no evidencias de sonificación EEG por sí mismas, ya que pueden producir sonido sin depender de la señal cerebral.

La matriz LED se mantiene como visualización lateral. Cuando está habilitada, reutiliza la información de notas recientes que también alimenta el piano roll. De este modo, LED y WebUI no constituyen motores musicales independientes, sino observadores de la misma actividad generada por el bloque de sonificación. Esta decisión mantiene coherencia entre la representación visual y la salida MIDI, sin introducir una segunda ruta de generación.

La Figura~\ref{fig:bloque_visualizacion_interaccion} resume la relación entre backend, servidor web, navegador y acciones auxiliares. La figura muestra que la WebUI recibe \textit{snapshots} del sistema, representa el estado en paneles y solo envía acciones acotadas de control musical o seguridad MIDI.

\begin{figure}[H]
    \centering
    \begin{tikzpicture}[
        webblock/.style={draw, rounded corners, align=center, text width=4.3cm, minimum height=0.95cm, font=\footnotesize},
        panelblock/.style={draw, rounded corners, align=center, text width=4.5cm, minimum height=0.95cm, font=\footnotesize, fill=gray!8},
        auxblock/.style={draw, rounded corners, dashed, align=center, text width=3.8cm, minimum height=0.9cm, font=\footnotesize, fill=gray!5},
        arrow/.style={-{Latex[length=3mm]}, thick},
        auxarrow/.style={-{Latex[length=2.5mm]}, dashed, thick},
        note/.style={align=center, font=\scriptsize}
    ]
        \node[webblock]   (backend) at (0,0)      {\texttt{BackendService}\\snapshot runtime};
        \node[webblock]   (server)  at (0,-1.6)   {\texttt{EEGWebServer}\\\texttt{GET /latest} + socket};
        \node[webblock]   (assets)  at (0,-3.2)   {assets WebUI\\HTML + JS + CSS};
        \node[webblock]   (browser) at (0,-4.8)   {navegador\\dashboard};
        \node[panelblock] (panels)  at (0,-6.4)   {paneles\\adquisición, calidad, sonificación, MIDI};

        \node[auxblock] (music) at (5.4,-3.2) {controles musicales\\root, main note y escala};
        \node[auxblock] (panic) at (5.4,-4.8) {\texttt{/midi/panic}\\seguridad MIDI};
        \node[auxblock] (roll)  at (5.4,-6.4) {piano roll / LED\\notas recientes};

        \draw[arrow] (backend) -- (server);
        \draw[arrow] (server) -- (assets);
        \draw[arrow] (assets) -- (browser);
        \draw[arrow] (browser) -- (panels);

        \draw[auxarrow] (music) -- (server);
        \draw[auxarrow] (panic) -- (server);
        \draw[auxarrow] (panels) -- (roll);

        \node[note] at (-4.45,-1.6) {La UI observa snapshots;\\no calcula DSP};
        \node[note] at (-4.45,-4.8) {No modifica ADS1299,\\firmware ni filtros MCU};
    \end{tikzpicture}
    \caption{Bloque de visualización e interacción. El backend publica snapshots hacia la WebUI, que representa el estado del sistema y permite acciones auxiliares acotadas: configuración musical, panic MIDI y visualización de notas recientes mediante piano roll o matriz LED.}
    \label{fig:bloque_visualizacion_interaccion}
\end{figure}

En resumen, el bloque de visualización e interacción se diseña como una capa auxiliar de observación, control musical limitado y seguridad operativa. Permite supervisar la adquisición, el procesamiento, la calidad de señal, la sonificación y el estado MIDI, pero no sustituye a la salida MIDI física ni participa en los cálculos críticos. Esta delimitación mantiene la arquitectura defendible: la WebUI ayuda a entender y controlar el sistema, mientras que la ruta principal sigue siendo EEG real, procesamiento, sonificación y MIDI OUT físico.

\subsection{Diseño electrónico y PCB}

% Pendiente de redacción propia.
% Describir conexiones principales: Arduino UNO Q, ADS1299, electrodos, salida MIDI, alimentación y consideraciones de seguridad.
% Fragmento incluido desde sections/diseno_sistema.tex mediante main.tex.
% La subsección \subsection{Diseño electrónico y PCB} se declara al final de diseno_sistema.tex.

\subsubsection{Antecedente del diseño electrónico base}

El diseño electrónico del prototipo parte de una arquitectura previa de adquisición de biopotenciales desarrollada en un trabajo anterior. Dicho antecedente planteaba una placa orientada a captar señales biomédicas de baja amplitud mediante un conversor analógico--digital específico, una etapa de protección y acondicionamiento de entrada, alimentación regulada, aislamiento galvánico y comunicación digital con un microcontrolador. Esta aproximación resulta adecuada como punto de partida, ya que separa de forma clara la zona conectada al usuario, la etapa de conversión analógica, el dominio digital y el sistema de comunicación.

En el presente TFG no se reproduce literalmente aquella arquitectura, sino que se toma como base conceptual para adaptarla a un objetivo diferente: transformar una señal EEG real en eventos musicales enviados mediante una salida MIDI física. Por tanto, el diseño electrónico no se entiende únicamente como una placa de adquisición, sino como el soporte físico de una cadena EEG--MIDI completa. Esta diferencia condiciona la elección del \textit{front-end}, la plataforma de control, la organización de la comunicación y la incorporación de una etapa final de salida MIDI.

La Figura~\ref{fig:adaptacion_diseno_base_pcb} resume esta transición. El diseño base se interpreta como una arquitectura de adquisición de biopotenciales, mientras que el diseño adaptado incorpora los elementos necesarios para cerrar la cadena hasta la salida musical física.

\begin{figure}[H]
    \centering
    \resizebox{0.95\textwidth}{!}{%
    \begin{tikzpicture}[
        block/.style={draw, rounded corners, align=center, text width=3.1cm, minimum height=0.95cm, font=\footnotesize},
        newblock/.style={draw, rounded corners, align=center, text width=3.1cm, minimum height=0.95cm, font=\footnotesize, fill=gray!8},
        arrow/.style={-{Latex[length=3mm]}, thick},
        note/.style={align=center, font=\scriptsize}
    ]
        \node[block] (base1) at (0,1.5) {Electrodos de biopotenciales};
        \node[block] (base2) at (3.8,1.5) {AFE biomédico del diseño base};
        \node[block] (base3) at (7.6,1.5) {Microcontrolador y comunicación};
        \node[block] (base4) at (11.4,1.5) {Datos de adquisición};

        \node[newblock] (new1) at (0,-1.5) {Electrodos EEG};
        \node[newblock] (new2) at (3.8,-1.5) {ADS1299--4PAGR, 4 canales y 24 bits};
        \node[newblock] (new3) at (7.6,-1.5) {Arduino UNO Q, MCU y Linux};
        \node[newblock] (new4) at (11.4,-1.5) {MIDI OUT físico};

        \draw[arrow] (base1) -- (base2);
        \draw[arrow] (base2) -- (base3);
        \draw[arrow] (base3) -- (base4);

        \draw[arrow] (new1) -- (new2);
        \draw[arrow] (new2) -- (new3);
        \draw[arrow] (new3) -- (new4);

        \node[note] at (5.7,0) {adaptación hacia sonificación EEG--MIDI};
        \draw[arrow] (5.7,0.85) -- (5.7,-0.85);
    \end{tikzpicture}}
    \caption{Adaptación conceptual del diseño electrónico base hacia la arquitectura EEG--MIDI. El objetivo deja de ser únicamente adquirir señales biomédicas y pasa a integrar adquisición EEG, procesamiento y salida MIDI física.}
    \label{fig:adaptacion_diseno_base_pcb}
\end{figure}

\subsubsection{Adaptación hacia ADS1299--4PAGR y Arduino UNO Q}

La primera adaptación relevante es la sustitución del conversor del diseño base por un ADS1299--4PAGR. Este circuito integrado pertenece a la familia ADS1299 de Texas Instruments, formada por conversores de bajo ruido, 24 bits y 4, 6 u 8 canales orientados a EEG y otras medidas de biopotenciales \cite{TexasInstrumentsADS1299}. En la placa desarrollada se emplea la variante de cuatro canales, coherente con el contrato final del sistema, donde el firmware y el backend Python mantienen una estructura fija de cuatro canales aunque la validación experimental se centre en CH1 como canal EEG principal.

La elección del ADS1299--4PAGR aporta varias ventajas para este prototipo. En primer lugar, integra entradas diferenciales de alta resolución, adecuadas para señales EEG de baja amplitud. En segundo lugar, incorpora funciones específicas de biopotenciales, como circuitería de BIAS, pines SRB, referencia y señales de diagnóstico. En tercer lugar, ofrece una interfaz SPI estándar, compatible con una lectura continua sincronizada mediante DRDY y adecuada para integrarse en una cadena embebida de adquisición en tiempo real.

La segunda adaptación importante es el cambio de plataforma de control hacia Arduino UNO Q. En esta arquitectura, la UNO Q cumple una doble función. El microcontrolador se encarga de la comunicación de bajo nivel con el ADS1299, de la temporización asociada a DRDY, del empaquetado inicial de las muestras y de la salida MIDI física. Paralelamente, el entorno Linux de la placa ejecuta el backend Python, que recibe los bloques EEG, calcula características, aplica la puerta de calidad y genera los eventos musicales. Esta separación permite mantener una adquisición determinista en el lado embebido y reservar las tareas de procesamiento y sonificación al entorno Python.

La adaptación también modifica la salida del sistema. El diseño base estaba orientado a adquisición y transmisión de datos, mientras que el sistema final incorpora una salida MIDI física. Esta decisión transforma la placa en un prototipo completo de interfaz EEG--MIDI: la señal no solo se observa o registra, sino que se convierte en mensajes musicales compatibles con dispositivos MIDI externos.

\subsubsection{Organización funcional de la placa}

La placa se organiza en bloques funcionales diferenciados. El primer bloque corresponde a las entradas de EEG y a la red de protección/acondicionamiento previa al conversor. El segundo bloque es el ADS1299--4PAGR, que actúa como núcleo de adquisición. El tercer bloque corresponde a la alimentación regulada y a la alimentación aislada. El cuarto bloque está formado por el aislamiento digital de las señales SPI y de control. El quinto bloque corresponde a la Arduino UNO Q, que conecta el dominio de adquisición con el firmware y el backend. Finalmente, el sexto bloque implementa la salida MIDI física mediante conector DIN-5 y etapa de adaptación.

Esta organización permite separar conceptualmente la zona conectada al usuario de la zona de control y de la salida musical. En los archivos KiCad de la placa aparecen, además, referencias de serigrafía que ayudan a interpretar esta división, como las zonas identificadas como \textit{Aislado} y \textit{No aislado}, la salida \textit{MIDI OUT}, los canales \texttt{ch1} a \texttt{ch4} y la señal \texttt{RLDRV}. Estas indicaciones no sustituyen a una validación formal de layout o seguridad eléctrica, pero sí muestran la intención de organizar físicamente la placa en dominios funcionales.

\begin{figure}[H]
    \centering
    \resizebox{0.98\textwidth}{!}{%
    \begin{tikzpicture}[
        block/.style={draw, rounded corners, align=center, text width=3.0cm, minimum height=0.95cm, font=\footnotesize},
        isol/.style={draw, rounded corners, align=center, text width=3.0cm, minimum height=0.95cm, font=\footnotesize, fill=gray!10},
        arrow/.style={-{Latex[length=3mm]}, thick},
        auxarrow/.style={-{Latex[length=2.5mm]}, dashed, thick},
        note/.style={align=center, font=\scriptsize}
    ]
        \node[block] (j1) at (0,0) {Conectores de electrodos};
        \node[block] (prot) at (3.7,0) {Protección y acondicionamiento};
        \node[block] (ads) at (7.4,0) {ADS1299--4PAGR, adquisición EEG};
        \node[isol] (adum) at (11.1,0) {ADuM3151, aislamiento SPI};
        \node[block] (uno) at (14.8,0) {Arduino UNO Q, MCU y Linux};
        \node[block] (midi) at (18.5,0) {DIN-5 MIDI OUT};

        \node[isol] (power) at (7.4,-2.3) {Alimentación aislada, CRE1S0505SC y reguladores};
        \node[block] (battery) at (0,-2.3) {Battery / 5 V, entrada de energía};

        \draw[arrow] (j1) -- (prot);
        \draw[arrow] (prot) -- (ads);
        \draw[arrow] (ads) -- (adum);
        \draw[arrow] (adum) -- (uno);
        \draw[arrow] (uno) -- (midi);

        \draw[auxarrow] (battery) -- (power);
        \draw[auxarrow] (power) -- (ads);
        \draw[auxarrow] (power) -- (adum);

        \node[note] at (7.4,1.35) {zona de adquisición conectada al usuario};
        \node[note] at (14.8,1.35) {zona de control y salida física};
    \end{tikzpicture}}
    \caption{Organización funcional de la placa. La PCB integra entradas EEG, protección, ADS1299--4PAGR, aislamiento digital, alimentación aislada, Arduino UNO Q y salida MIDI física.}
    \label{fig:organizacion_funcional_pcb}
\end{figure}

\subsubsection{Módulo de adquisición EEG}

El módulo de adquisición EEG está centrado en el ADS1299--4PAGR. Este componente constituye la frontera entre el dominio analógico de los electrodos y el dominio digital del sistema. En el lado analógico recibe los pares diferenciales asociados a los canales de EEG, así como las señales relacionadas con referencia, BIAS y alimentación analógica. En el lado digital expone las señales de comunicación y control necesarias para la lectura de datos: \texttt{DIN}, \texttt{DOUT}, \texttt{SCLK}, \texttt{CS}, \texttt{DRDY}, \texttt{START}, \texttt{RESET} y \texttt{PWDN}.

La variante de cuatro canales se ajusta al alcance experimental del prototipo. La arquitectura software conserva cuatro canales en el contrato de datos para mantener una estructura estable entre firmware, backend, capturas e interfaz auxiliar. Sin embargo, durante la validación final se trabaja con CH1 como canal principal de sonificación. Esta decisión no contradice el diseño multicanal, sino que acota la validación experimental a un canal EEG más controlado.

Desde el punto de vista de la memoria, este módulo debe describirse como el núcleo de adquisición y no como el bloque encargado de procesar la señal o generar música. El ADS1299 digitaliza la señal y entrega frames al firmware; el cálculo de bandas EEG, la evaluación de calidad y la generación de eventos musicales se realizan en etapas posteriores.

\subsubsection{Entradas, protección y referencia RLD/BIAS}

Las entradas de la placa se organizan mediante pares diferenciales asociados a los canales del ADS1299. En el diseño aparecen señales de entrada identificadas como \texttt{1P}/\texttt{1N}, \texttt{2P}/\texttt{2N}, \texttt{3P}/\texttt{3N} y \texttt{4P}/\texttt{4N}, junto con sus correspondientes señales internas hacia el conversor, como \texttt{IN\_1P}/\texttt{IN\_1N}, \texttt{IN\_2P}/\texttt{IN\_2N}, \texttt{IN\_3P}/\texttt{IN\_3N} e \texttt{IN\_4P}/\texttt{IN\_4N}. Esta nomenclatura facilita seguir el recorrido desde el conector de electrodos hasta el \textit{front-end}.

Entre los conectores y el ADS1299 se incorporan elementos de protección y acondicionamiento. La lista de materiales del diseño incluye dispositivos TPD4E1U06, redes resistivas de \SI{2.2}{\kilo\ohm} y redes capacitivas de \SI{1}{\nano\farad}. A nivel funcional, estos elementos permiten proteger las entradas frente a descargas electrostáticas y limitar la entrada de interferencias de alta frecuencia antes de la conversión analógico--digital. En este apartado se describen como parte del diseño de entrada, mientras que el cálculo detallado de frecuencias de corte o la validación de ruido queda reservado para análisis específicos.

La placa también contempla una red de referencia corporal asociada a las capacidades BIAS/RLD del ADS1299. En el esquemático aparecen los pines \texttt{BIASOUT}, \texttt{BIASIN}, \texttt{BIASINV} y \texttt{BIASREF}, así como señales identificadas como \texttt{RLD\_DRV} o \texttt{IN\_RLD\_DRV}. Esta red tiene como objetivo estabilizar el potencial común del montaje y mejorar el rechazo frente a interferencias de modo común. No obstante, al tratarse de un prototipo académico y no clínico, esta parte se describe como una solución experimental de referencia corporal, sin atribuirle prestaciones médicas ni certificación de seguridad.

\subsubsection{Alimentación e aislamiento galvánico}

El diseño de alimentación diferencia entre el lado de control y el lado asociado a la adquisición. La placa recibe alimentación desde la plataforma Arduino UNO Q y dispone de dominios de tensión utilizados por la electrónica digital y analógica. En la lista de materiales aparecen reguladores y conversores como \texttt{AMS1117--3.3}, \texttt{MIC5504--2.5YM5}, \texttt{TPS72325DBV} y \texttt{LM2664}. Estos componentes permiten generar las tensiones necesarias para el funcionamiento del ADS1299, su referencia y los circuitos auxiliares.

La alimentación aislada se apoya en el conversor DC--DC \texttt{CRE1S0505SC}. Este bloque conserva la filosofía del diseño base: separar la zona relacionada con el usuario de la zona de control del sistema. En el esquemático y la PCB aparecen dominios como \texttt{VDDA}, \texttt{VSSA}, \texttt{VDD} y \texttt{GNDD}, que ayudan a distinguir las referencias y alimentaciones utilizadas por las distintas partes del circuito.

La existencia de aislamiento no debe confundirse con una certificación clínica o con una validación normativa completa. En esta memoria se afirma únicamente que el diseño incorpora bloques de aislamiento y dominios diferenciados de alimentación. La comprobación de distancias, planos de masa, rigidez dieléctrica, fabricación y cumplimiento normativo pertenece a una validación electrónica posterior y queda fuera del alcance clínico del prototipo.

\subsubsection{Comunicación SPI aislada con Arduino UNO Q}

La comunicación entre el ADS1299 y la Arduino UNO Q se realiza mediante SPI, complementada con señales de control específicas del conversor. Las señales principales del bus son \texttt{MOSI}, \texttt{MISO}, \texttt{SCLK} y \texttt{CS}. A ellas se añaden \texttt{DRDY}, que indica la disponibilidad de una nueva conversión, y señales de control como \texttt{START}, \texttt{RESET} y \texttt{PWDN}.

En el diseño de la placa estas señales atraviesan un aislador digital \texttt{ADuM3151}. Por ello aparecen pares de señales antes y después de la barrera de aislamiento: \texttt{SPI1\_MOSI} se relaciona con \texttt{SPI\_MOSI\_ISO}, \texttt{SPI1\_MISO} con \texttt{SPI\_MISO\_ISO}, \texttt{SPI1\_SCLK} con \texttt{SPI\_SCLK\_ISO}, \texttt{FE\_SPI\_CS} con \texttt{SPI\_CS\_ISO}, \texttt{FE\_START} con \texttt{START\_ISO}, \texttt{FE\_RESET} con \texttt{RESET\_ISO} y \texttt{FE\_DRDY} con \texttt{DRDY\_ISO}. Esta estructura confirma que el ADS1299 no queda conectado directamente al microcontrolador, sino a través de una barrera de aislamiento digital.

\begin{figure}[H]
    \centering
    \small

    \setlength{\fboxsep}{6pt}

    \begin{tabular}{p{0.28\textwidth} c p{0.28\textwidth} c p{0.28\textwidth}}
        \centering
        \fbox{
        \begin{minipage}[t][1.8cm][c]{0.25\textwidth}
            \centering
            \textbf{Arduino UNO Q}\par
            Lado de control\par
            no aislado
        \end{minipage}
        }
        &
        \centering $\Longrightarrow$
        &
        \centering
        \fbox{
        \begin{minipage}[t][1.8cm][c]{0.25\textwidth}
            \centering
            \textbf{ADuM3151}\par
            Barrera de\par
            aislamiento digital
        \end{minipage}
        }
        &
        \centering $\Longrightarrow$
        &
        \centering
        \fbox{
        \begin{minipage}[t][1.8cm][c]{0.25\textwidth}
            \centering
            \textbf{ADS1299--4PAGR}\par
            Front-end EEG\par
            aislado
        \end{minipage}
        }
    \end{tabular}

    \vspace{0.5cm}

    \begin{tabular}{p{0.42\textwidth} c p{0.42\textwidth}}
        \fbox{
        \begin{minipage}[t]{0.39\textwidth}
            \textbf{SPI lado Arduino}\par
            \texttt{SPI1\_SCLK}\par
            \texttt{SPI1\_MOSI}\par
            \texttt{SPI1\_MISO}\par
            \texttt{FE\_SPI\_CS}
        \end{minipage}
        }
        &
        \centering $\Longleftrightarrow$
        &
        \fbox{
        \begin{minipage}[t]{0.39\textwidth}
            \textbf{SPI lado ADS1299}\par
            \texttt{SPI\_SCLK\_ISO}\par
            \texttt{SPI\_MOSI\_ISO}\par
            \texttt{SPI\_MISO\_ISO}\par
            \texttt{SPI\_CS\_ISO}
        \end{minipage}
        }
        \\[0.45cm]

        \fbox{
        \begin{minipage}[t]{0.39\textwidth}
            \textbf{Control y sincronismo}\par
            \texttt{FE\_START}\par
            \texttt{FE\_RESET}\par
            \texttt{FE\_DRDY}
        \end{minipage}
        }
        &
        \centering $\Longleftrightarrow$
        &
        \fbox{
        \begin{minipage}[t]{0.39\textwidth}
            \textbf{Control y sincronismo aislado}\par
            \texttt{START\_ISO}\par
            \texttt{RESET\_ISO}\par
            \texttt{DRDY\_ISO}
        \end{minipage}
        }
    \end{tabular}

    \vspace{0.45cm}

    \begin{minipage}{0.88\textwidth}
        \footnotesize
        La señal \texttt{DRDY} viaja desde el ADS1299 hacia la Arduino UNO Q e indica que existe una nueva conversión disponible. Las señales \texttt{START}, \texttt{RESET}, \texttt{CS}, \texttt{SCLK} y \texttt{MOSI} se gobiernan desde el lado de control.
    \end{minipage}

    \caption{Comunicación SPI aislada entre Arduino UNO Q y ADS1299--4PAGR. El ADuM3151 separa el lado de control del lado de adquisición, manteniendo las señales SPI, control y sincronismo necesarias para la lectura continua del conversor.}
    \label{fig:spi_aislado_adum3151}
\end{figure}

Esta comunicación aislada es coherente con el firmware, donde se definen las señales de selección de chip, reloj, entrada y salida SPI, DRDY, START, RESET y PWDN. El detalle de la secuencia de comandos del ADS1299, la escritura de registros, el modo RDATAC y la reconstrucción de frames pertenece al apartado de implementación. En el diseño electrónico basta con justificar la presencia de estas señales y su paso a través del aislamiento.

\subsubsection{Salida MIDI física}

La principal diferencia funcional del prototipo respecto a una placa de adquisición convencional es la presencia de una salida MIDI física. La PCB incluye un conector DIN-5 identificado como \texttt{MIDI OUT}. Esta salida se asocia a la señal \texttt{USART1\_TX} procedente de la Arduino UNO Q y a una etapa discreta formada por el transistor \texttt{MMBT5551} y resistencias de adaptación, entre ellas resistencias de \SI{220}{\ohm} y \SI{1}{\kilo\ohm}.

Desde el punto de vista de diseño del sistema, esta etapa permite que la cadena EEG--MIDI termine en una interfaz musical estándar. Los eventos generados por Python se convierten en bytes MIDI, se entregan al firmware mediante el contrato correspondiente y finalmente se escriben por la UART asociada a la salida física. La implementación validada utiliza \texttt{Serial1}/D1/TX y requiere inversión de polaridad de transmisión, resuelta en firmware. Por tanto, en este apartado se describe la etapa física de salida, mientras que la configuración UART, la inversión TX y el envío de mensajes \textit{Note On}/\textit{Note Off} se documentan en implementación.

La existencia de esta etapa es clave para delimitar el alcance del TFG. La WebUI y la matriz LED son herramientas auxiliares de observación, pero la salida principal del sistema es el conector MIDI físico. Esto permite conectar el prototipo a sintetizadores, módulos de sonido o interfaces MIDI externas, cumpliendo el objetivo de convertir actividad EEG en eventos musicales fuera del entorno de monitorización.

\subsubsection{Consideraciones de PCB y alcance del diseño}

La PCB integra físicamente los bloques principales descritos: Arduino UNO Q, ADS1299--4PAGR, protección de entradas, alimentación regulada, alimentación aislada, aislamiento digital, conectores de electrodos y salida MIDI física. La serigrafía y la asignación de redes permiten identificar la intención de separar zonas funcionales y de mantener una frontera entre adquisición, control y salida.

No obstante, la memoria debe limitar las afirmaciones al nivel verificable en los archivos de diseño disponibles. La identificación de componentes, conectores, señales y bloques funcionales permite describir la arquitectura de la placa, pero no equivale por sí sola a una validación completa del layout. No se deben sobreafirmar aspectos como cumplimiento normativo, distancias de aislamiento, fabricación definitiva, ruido real de la placa o comportamiento de los planos de masa si no han sido evaluados mediante ensayos o revisión específica.

En consecuencia, este subapartado presenta el diseño electrónico y PCB como la integración física de la arquitectura EEG--MIDI. Los detalles de configuración del ADS1299, lectura de frames, filtros en firmware, transporte por Bridge, procesamiento Python, sonificación y envío de bytes MIDI se desarrollan en los apartados de implementación. Del mismo modo, las capturas reales, pruebas de ruido, pruebas de BIAS/RLD, validación de CH1, benchmarks y comprobación de salida MIDI pertenecen al apartado de validación y resultados.



\end{document}
